\subsection{Simulation Framework}
\label{sub:sim_framework}

All of our numerical results are produced with the \texttt{sdim} qudit stabilizer simulator \cite{sdim}, which extends tableau-based stabilizer simulation to prime-dimensional qudits and samples noisy circuits in the Pauli frame picture. \autoref{sec:sdim-dem} reviews the simulator itself in more detail. The circuits in this work are expressed in terms of the generalized Hadamard, phase, $SUM$, and controlled-phase gates of \autoref{sec:bg}, computational and $X$-basis measurements, ancilla resets, and parameterized one- and two-qutrit Pauli noise channels.

For every combination of random Clifford sequence, benchmark depth, and stabilizer check schedule, we generate one annotated \texttt{sdim} circuit. The annotation includes the decomposed logical gates, the syndrome extraction blocks, noise placeholders positioned according to the chosen model of \autoref{sub:noise_models}, and \texttt{DETECTOR} and \texttt{LOGICAL\_OBSERVABLE} declarations that record which linear combinations of measurement outcomes are deterministic in the absence of noise. Each noise placeholder stores its rate as a multiplier relative to the swept parameter, so a single set of generated circuits serves the entire probability sweep. The logical Clifford catalog itself, along with the inverse lookup used to construct the recovery gate, is generated with a Cayley graph construction described in \autoref{sec:clifford-gen}.

The annotated circuits can be sampled through two backends. The first is direct frame simulation as described above. The second, which we use for the SI1000 sweep, compiles each circuit into a compact $q$-ary detector error model (DEM), a structure introduced for qubit stabilizer circuits by the Stim simulator \cite{Gidney2021Stim}. The DEM records how every possible fault in the circuit shifts each detector and logical observable modulo three, after which post-selected LRB-D statistics can be sampled at a cost that no longer depends on the size of the underlying stabilizer computation. \autoref{sec:sdim-dem} explains the construction, why it is exact for our noise models, and how it is sampled.
